\section{Experimental Methodology}
\label{sec:methodology}

\subsection{Two leaderboards that must not be merged}

\begin{description}
  \item[Provider benchmark.] Hosted services measured through their public
  APIs. The result includes the public network path, the control plane, the
  scheduler, the storage layer and any warm pool.
  \item[Infrastructure benchmark.] Firecracker, Cloud Hypervisor and gVisor
  deployed on identical hardware under our own
  control~\cite{firecracker2020,gvisor2019}.
\end{description}

The two are never ranked in one table. A hosted provider's CPU generation,
storage class, region and prewarming strategy are unknown and differ between
providers, so a hosted number and a controlled-hardware number do not measure
the same quantity.

\subsection{What every result publishes}

A sandbox benchmark measures a network path and an account configuration as
much as it measures a sandbox. Each result therefore carries: requested vCPU,
memory and disk; region or target; image or template identity; the client
machine and its network location; concurrency and arrival rate; warm or cold
state; whether persistence was enabled; the exec transport; and the wall-clock
window of the run. The framework records all of these automatically except the
client's network location, which the operator supplies.

\subsection{Statistics}

Each latency series is reported as $n$, median, $p_{90}$, $p_{95}$, $p_{99}$,
minimum and maximum, with a percentile bootstrap confidence interval on the
mean. We report medians rather than means as headline numbers: sandbox create
latency is long-tailed, and the mean of ten creates where one hit a cold pool
describes neither the warm case nor the cold one.

Repetitions are stated per benchmark. Runs are repeated on at least three
separate days, because provider performance varies with load that we cannot
observe, and a single day's measurement cannot be distinguished from that
day's conditions. \needsdata{the repetition schedule actually used}

\subsection{Cost}

Cost is reported per unit of completed work, not per hour:
\[
  C_{\mathrm{task}} = \frac{\text{total charge}}
  {\text{successfully completed tasks}},
\]
which charges startup overhead, execution time and failed retries to the work
they were spent on. For the reinforcement-learning profile the unit is a
valid rollout. \needsdata{prices, and the billing granularity each provider
actually applies}
